\section{Erd\H{o}s Problem \#284}
\subsection{Statement and attribution}
Let $f(k)$ be the largest possible smallest denominator in a representation
\[
 1=\frac1{n_1}+\cdots+\frac1{n_k},\qquad n_1<\cdots<n_k.
\]
Is $f(k)\sim k/(e-1)$? Yes. We reconstruct the deduction from Croot's short-interval theorem, paying attention to the fact that Croot allows the number of terms to vary, whereas here $k$ is prescribed. No novelty is claimed.

\subsection{Work I: an explicit upper bound}
Suppose $u=n_1>1$. Distinctness gives $n_i\ge u+i-1$, and the decreasing function $1/x$ gives
\[
 1\le\sum_{i=0}^{k-1}\frac1{u+i}
 \le\int_{u-1}^{u+k-1}\frac{dx}{x}
 =\log\left(1+\frac{k}{u-1}\right).
\]
Thus $u\le1+k/(e-1)$. For $u=1$ the same bound is immediate. Consequently
\[
 f(k)\le1+\frac{k}{e-1}.
 \tag{1}
\]

\subsection{External input and exact-length adjustment}
Croot's theorem, with target sum one, states that for every sufficiently large integer $N$ there is a finite set of distinct integers
\[
 N<b_1<\cdots<b_\ell<(e+\eta_N)N,
 \qquad \sum_{j=1}^{\ell}\frac1{b_j}=1,
 \qquad \eta_N\to0.
 \tag{2}
\]
This is the external number-theoretic input. It immediately bounds the number of terms by
$\ell\le(e-1+o(1))N$.

Any representation with largest denominator $b\ge2$ can be lengthened by exactly one using
\[
 \frac1b=\frac1{b+1}+\frac1{b(b+1)}.
\]
The new denominators are distinct and both exceed every retained denominator. Therefore repeated splitting of the current largest denominator preserves distinctness, preserves the reciprocal sum, and never decreases the smallest denominator.

Fix $0<\varepsilon<1$ and, for large $k$, take
$N=\lfloor(1-\varepsilon)k/(e-1)\rfloor$. For all sufficiently large $k$, (2) has $\ell\le k$. Apply the splitting operation $k-\ell$ times. The resulting representation has exactly $k$ terms and smallest denominator greater than $N$. Hence
\[
 \liminf_{k\to\infty}\frac{f(k)}k\ge\frac{1-\varepsilon}{e-1}.
\]
Letting $\varepsilon\downarrow0$ and combining with (1) proves $f(k)\sim k/(e-1)$.

\subsection{Verification and outcome}
There is no representation with two distinct terms, but representations exist for every $k\ge3$: start with $1=1/2+1/3+1/6$ and split the maximum. Thus the asymptotic function is well-defined for all large $k$; (1) makes its integer maximum finite. Splitting does not keep all denominators in Croot's interval, which is harmless because the problem controls only the smallest denominator.

\textbf{SOLVED, relative to the stated Croot theorem.} The upper bound and the passage from variable length to every prescribed large length are proved here.

\subsection{Sources}
E. S. Croot III, \emph{On unit fractions with denominators in short intervals}, Acta Arith. 99 (2001), 99--114, \url{https://arxiv.org/abs/math/9904181}; current statement: \url{https://www.erdosproblems.com/284}.
\section{COMPLETION ESTIMATE}
\textbf{COMPLETION: 100\%}
